\section{Teorema de Hahn-Banach}

\begin{ejercicio}
Sean $\varphi, \psi$ dos seminormas en un espacio vectorial $X$, y sea $f: X \longrightarrow \bb{K}$ un funcional lineal verificando que
\[
|f(x)| \leq \varphi(x) + \psi(x) \quad \forall x \in X
\]
Probar que existen dos funcionales lineales $g, h: X \longrightarrow \bb{K}$ tales que:
\[
f(x) = g(x) + h(x), \quad |g(x)| \leq \varphi(x), \quad |h(x)| \leq \psi(x) \quad \forall x \in X
\]
[Indicación: En el espacio vectorial $X \times X$ considérese la seminorma $p$ definida por $p(x_1, x_2) = \varphi(x_1) + \psi(x_2)$ y en el subespacio $\Delta = \{(x, x) : x \in X\}$ considérese el funcional lineal $h$ definido por $h(x, x) = f(x)$.]
\end{ejercicio}

\begin{ejercicio}
Sean $X$ un espacio normado real y $M$ un subespacio vectorial de $X$. Probar que si $f, g \in M^*$ verifican que $|f(m)| + |g(m)| \leq \|m\|$ para todo $m \in M$, entonces existen $h, k \in X^*$ que extienden a $f$ y $g$ respectivamente y verifican
\[
|h(x)| + |k(x)| \leq \|x\| \quad \forall x \in X
\]
[Indicación: Considérense $f + g$ y $f - g$.]
\end{ejercicio}

\begin{ejercicio}
Sea $X$ un espacio normado complejo, y $f: \bb{C} \longrightarrow X$ una función entera, es decir, derivable en todo punto de $\bb{C}$. Probar que, si $f$ está acotada, entonces $f$ es constante. (Se supone conocido el caso $X = \bb{C}$).
\end{ejercicio}

\begin{ejercicio}
Sea $X$ un espacio de Banach y $Z$ un subespacio cerrado de $X^*$, tal que $Z \neq X^*$. Supongamos que $Z$ separa los puntos de $X$, es decir,
\[
x \in X, \; x^*(x) = 0 \;\; \forall x^* \in Z \implies x = 0
\]
Probar que $X$ no es reflexivo.
\end{ejercicio}

\begin{ejercicio}
Probar que un subespacio de $c_0$, con la norma inducida por $c_0$, no puede ser isomorfo a $\ell_1$.
\end{ejercicio}

\begin{ejercicio}
Sean $X$ un espacio normado y $M$ un subespacio vectorial de $X$. Probar que para cada operador $T \in L(M, \ell_\infty)$, existe $S \in L(X, \ell_\infty)$ tal que $S|_M = T$ y $\|S\| = \|T\|$.
\end{ejercicio}

\begin{ejercicio}
Sean $X$ un espacio normado, $M$ un subespacio cerrado y $u \in X$. Probar que existe $f \in X^*$ tal que $|f(x)| \leq \mathrm{dist}(x, M)$ para todo $x \in X$ y $f(u) = \mathrm{dist}(u, M)$.
\end{ejercicio}

\begin{ejercicio}
Fijado $n \in \bb{N}$, probar que existe un funcional lineal y continuo $f$ en $C\left[0, 1\right]$ verificando que $f(p) = p'(0)$ para todo polinomio $p$ de grado menor o igual que $n$. ¿Existe un funcional lineal y continuo $f$ en $C\left[0, 1\right]$ tal que $f(p) = p'(0)$ para todo polinomio $p$?
\end{ejercicio}

\begin{ejercicio}
(Límites de Banach) Para cada $n \in \bb{N}$ consideramos $f_n: \ell_\infty \longrightarrow \bb{K}$ el funcional definido por
\[
f_n(x) = \frac{x(1) + x(2) + \dots + x(n)}{n} \quad (x \in \ell_\infty)
\]
Sea $M = \{x \in \ell_\infty : \{f_n(x)\} \text{ converge}\}$ y definamos $f(x) = \lim\limits_n f_n(x)$ ($x \in M$).
\begin{enumerate}[label=\textbf{\alph*)}]
    \item Probar que $f_n \in \ell_\infty^*$ y que $\|f_n\| = 1$ para todo $n \in \bb{N}$.
    \item Probar que $M$ es un subespacio vectorial de $\ell_\infty$ que contiene al espacio $c$ de las sucesiones convergentes.
    \item Probar que $f \in M^*$ con $\|f\| = 1$ y que $f(x) = \lim \{x_n\}$ para todo $x \in c$.
    \item Sea $\tau(x) = (x(2), x(3), \dots)$ para todo $x \in \ell_\infty$. Probar que $x - \tau(x) \in \ker(f) \subseteq M$ para todo $x \in \ell_\infty$.
    \item Deducir que existe $L \in \ell_\infty^*$, extensión de $f$, tal que $\|L\| = 1$ y $L(x) = L(\tau^n(x))$ para todo $x \in \ell_\infty$ y para todo $n \in \bb{N}$.
    \item Probar que $L(0, \nicefrac{1}{2}, 0, \nicefrac{1}{2}, \dots) = \nicefrac{1}{4}$. [Indicación: comparar con $L(\nicefrac{1}{2}, 0, \nicefrac{1}{2}, 0, \dots)$]. ¿Cuánto vale $L(1, 0, 0, 1, 0, 0, 1, \dots)$?
\end{enumerate}
\end{ejercicio}

\begin{ejercicio}
Sean $A$ y $B$ subconjuntos no vacíos, abiertos, convexos y disjuntos de un espacio normado $X$. Probar que existen $f \in X^*$ con $\|f\| = 1$ y $\alpha \in \bb{R}$ tales que $\mathrm{Re}\,f(a) < \alpha < \mathrm{Re}\,f(b)$ para todo $a \in A$ y $b \in B$. ¿Es siempre posible encontrar $f$ tal que $\sup \mathrm{Re}\,f(A) < \inf \mathrm{Re}\,f(B)$?
\end{ejercicio}